%Schriftgröße, Layout, Papierformat, Art des Dokumentes
\documentclass[footsepline,11pt,oneside,a4paper]{scrartcl}

\usepackage[left=2cm,right=2cm,top=2cm,bottom=2cm]{geometry}
\usepackage{scrpage2}
%mathem. funktionen
\usepackage{amsmath}
\usepackage{amsfonts}
\usepackage{amssymb}
%Umlaute ermöglichen
\usepackage[utf8]{inputenc}
%alte Rechtschreibung
\usepackage{german}
\usepackage{makeidx}
\usepackage{graphicx}
\usepackage[backend]{graphicx}
\usepackage{color}
\usepackage{units}
\usepackage{latexsym}
\usepackage{kpfonts}
\usepackage{fourier}
\author{Michael Meese / Thomas Fromme}
\title{Dokumentation Stereoskopie / StereoGL}
%\setheadwidth{1.5cm}

\newcommand{\hertz}[1]{\unit{#1}{Hz}}
\defpagestyle{GD}{%
  {\hfill}{\hfill}{\includegraphics[width=\textwidth]{img/fhlogoGD}}%
}{%
	{\pagemark\hfill}{\hfill\pagemark}{%
		Dokumentation: Grafische Datenvearbeitung	(M.Meese /T.Fromme)\hfill\pagemark}
}


\title{\hfill\\Stereoskopie in der grafischen Datenverarbeitung }
\author{%
	\begin{tabular}[c]{rl}
		Michael Meese & Mat. Nr. : 206 \\
		Thomas Fromme & Mat. Nr. : 206650 \\
	\end{tabular}
}
\date{%
	\begin{tabular}[c]{rl}
%		Versuchstag : &  \\
		Dokumentationsausarbeitung vom : & \today \\
	\end{tabular}
}

\begin{document}
\renewcommand{\titlepagestyle}{GD}
\maketitle
\pagestyle{GD}

\newpage
\newline
\tableofcontents


\newpage
\section{Geschichte und Konzept von stereoskopischen Bildern}
\subsection{Geschichte von stereoskopischen Bildern}
\newline
\newline

\newline
\newline
\newline
\newline

\subsection{Konzept von stereoskopischen Bildern}
\newline
\newline


\newline
\newline
\newline
\newline
\newline
\newline

\subsection{ }
\newline
\newline

\begin{tabbing}
\hspace{30}\=\kill
\textbullet \>  \\
\textbullet \>  \\
\textbullet \>  \\
\textbullet \>  \\
\textbullet \>  \\
\end{tabbing}
\newpage

\par
\subsection{Formeln}
\newline
\begin{center}
\large\underline{ }
\end{center}
\newline
\newline
\newline
\newline

\begin{displaymath}

\end{displaymath}
\newline
\newline

\newline
\newpage

\subsection{Implementation mit der OpenGL Bibliothek}
\newline



\newpage
\section{Diskussion und Vergleich beider Messverfahren}
\begin{center}
\textbf{Vergleich Oszilloskop und Timer / Counter }
\end{center}
\newline
\newline


\newline
\newpage
\section{Anhang:}

\newpage
\section{Quellen:}

\end{document}